\section{Variational Autoencoder (VAE)}

\subsection{从确定性到概率性编码}

传统 Autoencoder 的编码是确定性的：给定输入 $x$，编码器输出固定的隐向量 $z$。这使得隐空间可能存在"空洞"和不连续区域，从这些区域采样将产生无意义的输出。

\begin{figure}[H]
\centering
\includegraphics[width=0.85\textwidth]{slides-images/slide-20.jpg}
\caption{传统 Autoencoder 的确定性编码流程\protect\footnotemark}
\end{figure}
\footnotetext{Slides 第20页。}

\textbf{Variational Autoencoder (VAE)} 引入了概率性的改进：编码器不再输出固定的隐向量，而是输出一个\textbf{概率分布的参数}——均值向量 $\mu$ 和标准差向量 $\sigma$。然后从该分布中\textbf{采样}得到隐向量 $z$。

\begin{figure}[H]
\centering
\includegraphics[width=0.85\textwidth]{slides-images/slide-22.jpg}
\caption{VAE 的核心改进：概率性隐空间\protect\footnotemark}
\end{figure}
\footnotetext{Slides 第22页。来源：Kingma+ ICLR 2014。}

\begin{importantbox}{VAE 的核心思想}
VAE 是 Autoencoder 的\textbf{概率性改进}。编码器输出均值 $\mu$ 和标准差 $\sigma$，定义一个关于隐变量的概率分布 $q_\phi(z|x)$。通过从该分布中采样，VAE 可以生成不同的隐向量，从而实现多样化的生成。
\end{importantbox}

\subsection{VAE 的损失函数}

VAE 的优化目标由两部分组成：

\begin{figure}[H]
\centering
\includegraphics[width=0.85\textwidth]{slides-images/slide-25.jpg}
\caption{VAE 优化：重建损失 + 正则化项\protect\footnotemark}
\end{figure}
\footnotetext{Slides 第25页。来源：Kingma+ ICLR 2014。}

$$\mathcal{L}(\phi, \theta, x) = \underbrace{\mathbb{E}_{q_\phi(z|x)}[\log p_\theta(x|z)]}_{\text{重建损失（Reconstruction Loss）}} - \underbrace{D_{KL}(q_\phi(z|x) \| p(z))}_{\text{正则化项（Regularization Term）}}$$

\begin{itemize}
    \item \textbf{重建损失}：衡量解码器从隐变量 $z$ 重建数据 $\hat{x}$ 的质量，鼓励 $\hat{x}$ 尽可能接近 $x$
    \item \textbf{正则化项（KL Divergence）}：衡量编码器输出的隐变量分布 $q_\phi(z|x)$ 与先验分布 $p(z)$ 之间的距离，鼓励隐空间具有良好的结构
    \item $\phi$：编码器的参数
    \item $\theta$：解码器的参数
\end{itemize}

\subsection{先验分布与 KL Divergence}

VAE 需要对隐变量设定一个\textbf{先验分布}（Prior）$p(z)$。最常见的选择是标准正态分布：

$$p(z) = \mathcal{N}(\mu = 0, \sigma^2 = 1)$$

\begin{figure}[H]
\centering
\includegraphics[width=0.85\textwidth]{slides-images/slide-28.jpg}
\caption{先验分布：正态分布鼓励隐变量均匀分布在隐空间中心\protect\footnotemark}
\end{figure}
\footnotetext{Slides 第28页。来源：Kingma+ ICLR 2014。}

\begin{knowledgebox}{为什么选择正态先验？}
选择 $\mathcal{N}(0, 1)$ 作为先验有两个重要目的：
\begin{enumerate}
    \item \textbf{集中编码}：鼓励编码器将隐变量分布在隐空间的中心区域
    \item \textbf{防止"作弊"}：防止网络将不同类别的数据点聚集到隐空间的特定孤立区域（即记忆数据），而是迫使它学习有意义的、连续的表示
\end{enumerate}
\end{knowledgebox}

\textbf{KL Divergence}（KL 散度）是衡量两个概率分布之间"距离"的指标。对于正态分布先验，KL 散度有解析表达式：

$$D_{KL}(q_\phi(z|x) \| p(z)) = -\frac{1}{2} \sum_{j=1}^{J} \left(1 + \log(\sigma_j^2) - \mu_j^2 - \sigma_j^2\right)$$

其中 $J$ 是隐变量的维度，$\mu_j$ 和 $\sigma_j$ 分别是第 $j$ 个隐变量的均值和标准差。

\subsection{正则化的直觉：连续性与完整性}

正则化的作用是确保隐空间具有两个关键性质：

\begin{figure}[H]
\centering
\includegraphics[width=0.85\textwidth]{slides-images/slide-30.jpg}
\caption{正则化确保隐空间的连续性和完整性\protect\footnotemark}
\end{figure}
\footnotetext{Slides 第30页。来源：Joseph Rocca。}

\begin{enumerate}
    \item \textbf{连续性（Continuity）}：在隐空间中相邻的点，解码后应产生语义相似的内容
    \item \textbf{完整性（Completeness）}：从隐空间的任意位置采样，都应解码出有意义的内容
\end{enumerate}

\begin{figure}[H]
\centering
\includegraphics[width=0.85\textwidth]{slides-images/slide-32.jpg}
\caption{正则化后隐空间的连续性与完整性得到保证\protect\footnotemark}
\end{figure}
\footnotetext{Slides 第32页。来源：Joseph Rocca。}

\begin{warningbox}{无正则化的隐空间问题}
如果没有正则化（仅使用重建损失），网络可能学到的隐空间存在大量"空洞"——这些区域没有被训练数据覆盖。从空洞区域采样将产生无意义的噪声输出。正则化通过迫使隐变量分布接近标准正态分布，填充了这些空洞。
\end{warningbox}

\subsection{Reparameterization Trick}

VAE 训练面临一个技术难题：隐变量 $z$ 是通过\textbf{采样}操作得到的，而采样是不可微的——无法通过采样层进行反向传播。

\begin{figure}[H]
\centering
\includegraphics[width=0.85\textwidth]{slides-images/slide-35.jpg}
\caption{Reparameterization Trick：使采样可微\protect\footnotemark}
\end{figure}
\footnotetext{Slides 第35页。来源：Kingma+ ICLR 2014。}

解决方案是\textbf{Reparameterization Trick}（重参数化技巧）。核心思想是将随机性从隐变量 $z$ 中分离出来：

$$z = \mu + \sigma \odot \epsilon, \quad \epsilon \sim \mathcal{N}(0, 1)$$

\begin{itemize}
    \item $\mu$：编码器输出的均值向量（可微的确定性节点）
    \item $\sigma$：编码器输出的标准差向量（可微的确定性节点）
    \item $\epsilon$：从标准正态分布采样的随机噪声（固定的随机常数）
    \item $\odot$：逐元素乘法
\end{itemize}

\begin{importantbox}{Reparameterization Trick 的精髓}
通过将 $z$ 重写为 $\mu$、$\sigma$ 的确定性函数加上一个外部随机噪声 $\epsilon$，梯度可以通过 $\mu$ 和 $\sigma$ 正常反向传播。随机性被"外部化"为从固定分布采样的常数，不阻断梯度流。这使得整个 VAE 可以端到端地用反向传播训练。
\end{importantbox}

\subsection{隐变量扰动与可解释性}

训练好的 VAE 具有一个重要特性：隐空间中的不同维度往往对应不同的\textbf{可解释特征}。

\begin{figure}[H]
\centering
\includegraphics[width=0.85\textwidth]{slides-images/slide-38.jpg}
\caption{隐变量扰动：固定其他变量，改变单个变量观察语义变化\protect\footnotemark}
\end{figure}
\footnotetext{Slides 第38页。来源：Kingma+ ICLR 2014。}

例如，在人脸数据集上训练的 VAE 中，固定所有隐变量但缓慢改变其中一个维度的值，可以观察到生成的人脸在某个特定属性（如头部姿态、表情等）上的平滑变化。这说明 VAE 成功地学习到了数据中隐含的语义特征。

\subsection{$\beta$-VAE 与隐空间解耦}

标准 VAE 的隐空间中，不同维度可能编码了纠缠在一起的特征。$\beta$-VAE 通过在正则化项前添加一个超参数 $\beta > 1$ 来加强解耦：

$$\mathcal{L}(\theta, \phi; x, z, \beta) = \mathbb{E}_{q_\phi(z|x)}[\log p_\theta(x|z)] - \beta \cdot D_{KL}(q_\phi(z|x) \| p(z))$$

\begin{figure}[H]
\centering
\includegraphics[width=0.85\textwidth]{slides-images/slide-40.jpg}
\caption{$\beta$-VAE 实现更好的隐空间解耦\protect\footnotemark}
\end{figure}
\footnotetext{Slides 第40页。来源：Higgins+ ICLR 2017。}

当 $\beta = 1$ 时退化为标准 VAE；$\beta > 1$ 时更强地约束隐空间瓶颈，鼓励每个隐变量独立编码不同的语义特征（disentanglement），从而改变一个属性时不会影响其他属性。

\subsection{VAE 总结}

\begin{figure}[H]
\centering
\includegraphics[width=0.85\textwidth]{slides-images/slide-45.jpg}
\caption{VAE 核心要点总结\protect\footnotemark}
\end{figure}
\footnotetext{Slides 第45页。}

VAE 的五个核心要点：
\begin{enumerate}
    \item \textbf{压缩表示}：将世界的高维表示压缩为可学习的低维隐空间
    \item \textbf{无监督学习}：通过重建损失实现无标签学习
    \item \textbf{Reparameterization Trick}：使含采样层的网络可端到端训练
    \item \textbf{可解释隐变量}：通过扰动分析理解各隐变量对应的语义特征
    \item \textbf{生成新样本}：从隐空间采样并解码即可生成新数据
\end{enumerate}

\subsection{本章小结}

VAE 是 Autoencoder 的概率性扩展，通过引入概率隐空间和 KL Divergence 正则化，解决了传统 Autoencoder 隐空间不连续、不完整的问题。Reparameterization Trick 巧妙地解决了通过采样层反向传播的难题。VAE 是第一个被介绍的深度生成模型，它属于 Latent Variable Model 家族。


